\documentclass{beamer}
\usepackage[latin1]{inputenc}
\beamertemplateshadingbackground{red!10}{blue!10}
%\usepackage{fancybox}
\usepackage{epsfig}
\usepackage{verbatim}
\usepackage{url}
%\usepackage{graphics}
%\usepackage{xcolor}
\usepackage{fancybox}
\usepackage{moreverb}
%\usepackage[all]{xy}
\usepackage{listings}
\usepackage{filecontents}
\usepackage{graphicx}

\lstset{
  language=Lisp,
  basicstyle=\scriptsize\ttfamily,
  keywordstyle={},
  commentstyle={},
  stringstyle={}}

\def\inputfig#1{\input #1}
\def\inputeps#1{\includegraphics{#1}}
\def\inputtex#1{\input #1}

\inputtex{logos.tex}

%\definecolor{ORANGE}{named}{Orange}

\definecolor{GREEN}{rgb}{0,0.8,0}
\definecolor{YELLOW}{rgb}{1,1,0}
\definecolor{ORANGE}{rgb}{1,0.647,0}
\definecolor{PURPLE}{rgb}{0.627,0.126,0.941}
\definecolor{PURPLE}{named}{purple}
\definecolor{PINK}{rgb}{1,0.412,0.706}
\definecolor{WHEAT}{rgb}{1,0.8,0.6}
\definecolor{BLUE}{rgb}{0,0,1}
\definecolor{GRAY}{named}{gray}
\definecolor{CYAN}{named}{cyan}

\newcommand{\orchid}[1]{\textcolor{Orchid}{#1}}
\newcommand{\defun}[1]{\orchid{#1}}

\newcommand{\BROWN}[1]{\textcolor{BROWN}{#1}}
\newcommand{\RED}[1]{\textcolor{red}{#1}}
\newcommand{\YELLOW}[1]{\textcolor{YELLOW}{#1}}
\newcommand{\PINK}[1]{\textcolor{PINK}{#1}}
\newcommand{\WHEAT}[1]{\textcolor{wheat}{#1}}
\newcommand{\GREEN}[1]{\textcolor{GREEN}{#1}}
\newcommand{\PURPLE}[1]{\textcolor{PURPLE}{#1}}
\newcommand{\BLACK}[1]{\textcolor{black}{#1}}
\newcommand{\WHITE}[1]{\textcolor{WHITE}{#1}}
\newcommand{\MAGENTA}[1]{\textcolor{MAGENTA}{#1}}
\newcommand{\ORANGE}[1]{\textcolor{ORANGE}{#1}}
\newcommand{\BLUE}[1]{\textcolor{BLUE}{#1}}
\newcommand{\GRAY}[1]{\textcolor{gray}{#1}}
\newcommand{\CYAN}[1]{\textcolor{cyan }{#1}}

\newcommand{\reference}[2]{\textcolor{PINK}{[#1~#2]}}
%\newcommand{\vect}[1]{\stackrel{\rightarrow}{#1}}

% Use some nice templates
\beamertemplatetransparentcovereddynamic

\newcommand{\A}{{\mathbb A}}
\newcommand{\degr}{\mathrm{deg}}

\title{First-class global environments for \commonlisp{}}

\author{Robert Strandh}
\institute{
LaBRI, University of Bordeaux
}
\date{April, 2015}

%\inputtex{macros.tex}


\begin{document}
\frame{
\resizebox{3cm}{!}{\includegraphics{Logobx.pdf}}
\hfill
\resizebox{1.5cm}{!}{\includegraphics{labri-logo.pdf}}
\titlepage
\vfill
\small{European Lisp Symposium, London, UK \hfill ELS2015}
}

\setbeamertemplate{footline}{
\vspace{-1em}
\hspace*{1ex}{~} \GRAY{\insertframenumber/\inserttotalframenumber}
}

\frame{
\frametitle{Context: The \sicl{} project}

https://github.com/robert-strandh/SICL

Several objectives:

\begin{itemize}
\item Create high-quality \emph{modules} for implementors of
  \commonlisp{} systems.
\item Improve existing techniques with respect to algorithms and data
  structures where possible.
\item Improve readability and maintainability of code.
\item Improve documentation.
\item Ultimately, create a new implementation based on these modules.
\end{itemize}
}

\frame{
\frametitle{Environments in \commonlisp{}}

}

\frame{
\frametitle{Environments in the \hs{}}

Mostly used to distinguish the \emph{global environment} from a
\emph{non-local} (lexical) environment.
\vskip 0.25cm
For example, \texttt{setf} might need to call \texttt{macroexpand}
with the same environment, in case it is wrapped in \texttt{macrolet}
or \texttt{symbol-macrolet}.
}

\frame{
\frametitle{Global environments in the \hs{}}

Occasionally, used to distinguish between the compile-time
\emph{evaluation environment} and the global run-time environment.
\vskip 0.25cm
Example from \texttt{find-class}:
\begin{quotation}
  The \emph{environment} might be used to distinguish between a
  compile-time and a run-time environment.
\end{quotation}
\vskip 0.25cm
Other examples: \texttt{typep}, \texttt{subtypep}, \texttt{constantp},
and \texttt{compiler-macro-function}.
\vskip 0.25cm
However, the \hs{} explicitly allows for all global environments to be
one and the same.

}

\frame{
\frametitle{Global environments in typical implementations}

A typical implementation has a single global environment as allowed by
the \hs{}.  It is typically \emph{spread out} in that a \emph{symbol}
might contain a \emph{function cell} and/or a \emph{value cell}.
\vskip 0.25cm
Some implementations store the function named \emph{foo} in a slot of
the symbol named \emph{foo} but they store the function named
\texttt{(setf \textsf{\emph{foo}})} elsewhere.
\vskip 0.25cm
This discrepancy might reflect the fact that the basic implementation
pre-dates the introduction of \texttt{setf} functions.

}

\frame{
\frametitle{Previous work}
}

\frame{
\frametitle{First-class global environments}

In \sicl{}, we put all environment-related information in an
\texttt{environment} class instance.  This environment would typically
be organized as a collection of hash tables, one for each
\emph{namespace}:
\begin{itemize}
\item The \emph{function} namespace, containing information about
  functions, macros, and special operators.
\item The \emph{variable} namespace, containing information about
  special variables, constant variables, and symbol macros.
\item The \emph{type} namespace, containing types defined by
  \texttt{deftype}.
\item The \emph{class} namespace, containing a mapping from class
  names (symbols) to class metaobjects.
\item The \emph{package} namespace, containing mappings from package
  names (strings) to package objects.
\item etc.
\end{itemize}
}

\frame{
\frametitle{A protocol for first-class global environments}

However, the current paper is only about a \clos{}-based
\emph{protocol} for first-class global environments.

}


\frame{
  \frametitle{Acknowledgments}

We would like to thank Alastair Bridgewater, David Murray, Robert
Smith, Nicolas Hafner, and Bart Botta for providing valuable feedback
on early versions of this paper.
}

\frame{
\frametitle{Future work}

\begin{itemize}
\item Extend protocol to include environment \emph{inheritance}.
\item Design technique for efficient environment cloning.
\item Create a multi-user system based on first-class global
  environments.
\end{itemize}

}

\frame{
\frametitle{Conclusion and Perspectives}
\vfill
\hfill
\ORANGE{Thank you for your attention!}
}

%% \frame{\tableofcontents}
%% \bibliography{references}
%% \bibliographystyle{alpha}

\end{document}
